\chapter*[Introdução]{Introdução}
\addcontentsline{toc}{chapter}{Introdução}

O presente trabalho tem como objetivo propor uma arquitetura baseada em blockchain para o gerenciamento
de permissões de acesso a prontuários eletrônicos de pacientes. A proposta busca permitir o compartilhamento
de informações médicas entre diferentes instituições de saúde de forma segura e auditável, respeitando os princípios
estabelecidos pela Lei Geral de Proteção de Dados (LGPD).

Nesse contexto, considera-se que os dados de prontuários não podem ser distribuídos entre hospitais sem o consentimento
do paciente, uma vez que se tratam de dados pessoais sensíveis. Assim, o sistema proposto utiliza contratos inteligentes
para registrar e gerenciar o consentimento do titular dos dados, permitindo que o paciente autorize ou restrinja o
acesso às suas informações médicas.

Adicionalmente, o modelo também contempla situações excepcionais previstas em contextos clínicos, como emergências médicas,
nas quais o acesso aos dados pode ser necessário mesmo na ausência de autorização prévia do paciente.
Dessa forma, a arquitetura busca equilibrar a proteção da privacidade, a conformidade com a legislação
vigente e a disponibilidade de informações essenciais para o atendimento em saúde.
